% Propozycja zastępująca content/04-metodyka.tex.
% Nie jest włączona do main.tex. Przed włączeniem rozstrzygnąć uwagi poniżej.
% Podstawa: instrukcja J. Piechowicza, lokalny konspekt, s. 2--4.
% Dodać rzeczywisty wpis instrukcji do references.bib oraz odwołanie bibliograficzne.
\section{Metodyka}
\label{sec:metodyka}

\subsection{Sposób wykonania pomiarów}
Pomiary wykonano metodą próbkowania opisaną w instrukcji
laboratoryjnej. Przed rozpoczęciem badań przeprowadzono oględziny
obiektu i jego otoczenia oraz wybrano punkt pomiaru hałasu
oraz punkt pomiaru tła akustycznego, przedstawione na
rysunku~\ref{fig:punkty_pomiarowe}.
Mikrofon w punkcie pomiaru hałasu umieszczono na wysokości
$1{,}5\,\mathrm{m}$ nad gruntem, po północnej stronie stacji.

Rejestrowano równoważny poziom dźwięku A w elementarnych
przedziałach pomiarowych o czasie $t_0=10\,\mathrm{s}$.
Zapisane pliki miernika SVAN~971 potwierdzają zastosowanie
charakterystyki częstotliwościowej A i detektora Fast w pierwszym
profilu pomiarowym oraz liniowego całkowania dla wyniku $L_{Aeq}$.
Każdy plik zawiera wynik dla całego przedziału $10\,\mathrm{s}$
oraz dziesięć rekordów rejestracji z krokiem $1\,\mathrm{s}$.
Rekordów tych nie traktowano jako odrębnych próbek
$10\,\mathrm{s}$.

% KONIECZNE: potwierdzić przyczynę skrócenia t0 do 10 s.
% Jeśli obserwacje to potwierdzają, dopisać:
% Czas próbki skrócono do 10 s ze względu na przejazdy pojazdów,
% pomiędzy którymi nie było możliwe wykonanie niezakłóconej próbki 60 s.
% Bez potwierdzenia nie wpisywać tego zdania jako faktu.
Zgodnie z instrukcją zastosowanie próbek $10\,\mathrm{s}$
zamiast $60\,\mathrm{s}$ wymaga występowania okresowych zakłóceń
uniemożliwiających wykonanie dłuższej próbki przy stałym hałasie
badanego źródła. Uzasadnienie tego wyboru na podstawie obserwacji
terenowych wymaga uzupełnienia dokumentacji ćwiczenia.

\subsection{Kontrola liczby próbek}
Dla każdej serii wyznacza się rozstęp wyników
\begin{equation}
 R=\max_k L_{Ak}-\min_k L_{Ak},
\end{equation}
gdzie $L_{Ak}$ oznacza poziom równoważny zmierzony w $k$-tej
próbce o czasie $t_0$.
Liczbę wykonanych próbek porównuje się z wymaganiami instrukcji:
3 próbki dla $0<R\leq1\,\mathrm{dB}$,
4 dla $1<R\leq1{,}5\,\mathrm{dB}$,
5 dla $1{,}5<R\leq2\,\mathrm{dB}$,
6 dla $2<R\leq2{,}5\,\mathrm{dB}$ oraz
7 dla $2{,}5<R\leq3\,\mathrm{dB}$.
Przy rozstępie większym niż $3\,\mathrm{dB}$ instrukcja
przewiduje wydłużenie czasu próbki.
Wyniki sprawdzenia liczebności serii należy przedstawić
wraz z wynikami pomiarów, również w przypadku niespełnienia
tego wymagania.

\subsection{Wyznaczanie tła akustycznego}
Ze względu na brak możliwości wyłączenia badanego obiektu
pomiar tła wykonano po pomiarze zasadniczym, w odrębnym punkcie
zaznaczonym na rysunku~\ref{fig:punkty_pomiarowe}.
Tło powinno reprezentować dźwięki niezwiązane z badaną stacją,
przy porównywalnym oddziaływaniu pozostałych źródeł.
% UZUPEŁNIĆ na podstawie obserwacji: dlaczego punkt tła był
% porównywalny (ulice, przesłonięcie stacji, odległości, pora, pogoda).
% Potwierdzić przypisanie P04--P08 oraz P10--P14 przed dalszym opracowaniem.
Ocena porównywalności punktów oraz przypisanie plików do serii
zasadniczej i serii tła wymagają potwierdzenia na podstawie
notatek terenowych.

\subsection{Opracowanie wyników}
Dla próbek o jednakowym czasie trwania średni poziom
w serii wyznacza się przez uśrednienie energetyczne:
\begin{equation}
 L_{A\mathrm{sr}}=10\log_{10}\left(
 \frac{1}{n}\sum_{k=1}^{n}10^{0{,}1L_{Ak}}\right).
 \label{eq:srednia-energetyczna}
\end{equation}
Analogicznie wyznacza się średni poziom tła $L_{At}$
z próbek przypisanych do pomiaru tła.
Poziom związany z pracą badanego obiektu w rozpatrywanym
przedziale otrzymuje się przez odejmowanie energetyczne:
\begin{equation}
 L_{Aek}=10\log_{10}\left(
 10^{0{,}1L_{A\mathrm{sr}}}-10^{0{,}1L_{At}}\right).
 \label{eq:korekcja-tla}
\end{equation}
Zgodnie z instrukcją przy różnicy
$L_{A\mathrm{sr}}-L_{At}<3\,\mathrm{dB}$
poziom źródła wymaga wyznaczenia metodą obliczeniową;
powyższego odejmowania nie traktuje się wówczas jako
wystarczającego wyznaczenia wyniku.
Przy różnicy większej niż $10\,\mathrm{dB}$ instrukcja
dopuszcza pominięcie wpływu tła.

Wynik uzyskany z krótkiej serii charakteryzuje badane warunki
pracy i punkt pomiarowy. Odniesienie go do wskaźnika dla pory
dnia lub nocy wymaga dodatkowej informacji o czasie pracy
źródeł i reprezentatywności zarejestrowanej serii.
